\section{Conclusion}\label{sec:conclusion}
\subsection{Multi Criteria Analysis}
The multi-criteria analysis shows that all three frameworks meet the essential communication requirements, as they all support both RESTful API calls and WebSocket connections.
Differences between the frameworks become more apparent in the weighted criteria such as Community, Developer Experience, and Ecosystem.

Spring Boot scores slightly higher than Echo in the overall evaluation, mainly due to a larger community and a more mature ecosystem.
NestJS, however, achieves the highest final score of 8.38 out of 10, largely driven by its strong ecosystem and developer experience.
This indicates that while Spring Boot is well-established and widely supported, NestJS provides a more flexible and modern development experience.

Quality Assurance scores are equal across all frameworks, reflecting that each provides sufficient tools and support to build secure and reliable applications.
Developer experience and ecosystem differences are therefore the primary factors influencing the final ranking.

Overall, the analysis suggests that NestJS is the most suitable choice for the project based on the selected criteria, while Spring Boot remains a close alternative, especially for teams with prior Java experience.
Echo, while capable, ranks lower due to a smaller community and ecosystem support.

\subsection{Prototyping}
Based on the prototyping experience, both frameworks are capable and well-documented, but the choice depends on priorities.

Spring Boot offers a mature, stable environment with strong IDE integration and a large community.
Its annotation-based approach and built-in features make it easy to set up controllers, services, repositories, and testing.
Observability and security can be configured reliably, although some examples are out of date for newer versions.

NestJS provides a modern, TypeScript-based environment with a modular architecture and dependency injection.
Setting up controllers, services, and repositories is straightforward, and the framework integrates well with ORMs like TypeORM.
Some initial challenges, such as validation and handling complex entity relations, required workarounds, but documentation was clear and helpful.

Considering developer experience, ecosystem, and flexibility, NestJS emerges as the better choice for this project.
